% ============================================================
%  学术论文/报告模板 —— LaTeX article + 常用宏包
%  编译方式: XeLaTeX（中文推荐），论文需编译两次生成引用
%  用法: 替换【】内容，参考文献放入 refs.bib
% ============================================================
\documentclass[12pt,a4paper]{article}

% ---------- 页面 ----------
\usepackage[margin=2.5cm]{geometry}

% ---------- 中文 ----------
\usepackage{ctex}

% ---------- 数学 ----------
\usepackage{amsmath,amssymb,amsfonts}
\usepackage{bm}                    % 粗体数学符号
\newcommand{\T}{^{\mathsf{T}}}     % 转置
\newcommand{\R}{\mathbb{R}}        % 实数集
\newcommand{\mat}[1]{\bm{#1}}      % 向量/矩阵粗体

% ---------- 图表 ----------
\usepackage{graphicx}
\usepackage{booktabs}              % 三线表
\usepackage{float}                 % [H] 强制位置
\usepackage{caption}

% ---------- 引用与链接 ----------
\usepackage[numbers,sort&compress]{natbib}
\usepackage{hyperref}
\hypersetup{colorlinks=true, citecolor=blue, linkcolor=black}

% ---------- 代码 ----------
\usepackage{listings}
\usepackage{xcolor}
\lstset{
  basicstyle=\ttfamily\small,
  numbers=left,
  numberstyle=\tiny,
  frame=single,
  breaklines=true,
  language=Python,
  commentstyle=\color{gray},
  keywordstyle=\color{blue}
}

% ============================================================
%	论文元信息
% ============================================================
\title{\textbf{论文题目：研究的核心贡献}}
\author{作者一$^{1}$，作者二$^{2}$ \\[0.3em]
\small $^{1}$单位一，$^{2}$单位二}
\date{\today}

\begin{document}
\maketitle

% ---------- 摘要 ----------
\begin{abstract}
【200-300 字：研究背景一句话 → 本文解决的问题 → 提出的方法 → 关键结果（含数值） → 意义。】
\par\medskip
\textbf{关键词：}关键词一；关键词二；关键词三
\end{abstract}

% ============================================================
\section{引言 (Introduction)}
【研究背景与意义：由大到小引出问题。】
\begin{enumerate}
    \item 领域重要性及现有进展；
    \item 现有方法的三点不足（可用引用支撑，如 \cite{ref1}）；
    \item 本文方法的核心思想与创新点（1-2-3 列出）；
    \item 主要贡献总结 + 全文组织结构。
\end{enumerate}

% ============================================================
\section{相关工作 (Related Work)}
【分主题梳理文献，重点放在与本文的差异上，避免罗列式堆砌。】
\subsection{方向一}
\subsection{方向二}

% ============================================================
\section{问题描述与数学建模 (Problem Formulation)}
\begin{equation}
    \min_{\mat{\theta}} \; \mathcal{L}(\mat{\theta}) = \frac{1}{N}\sum_{i=1}^{N} \ell\big(f(\mat{x}_i; \mat{\theta}), y_i\big) + \lambda \|\mat{\theta}\|_2^2
    \label{eq:objective}
\end{equation}
其中 $\mat{x}_i \in \R^d$ 表示第 $i$ 个样本，$y_i$ 为标签，$f(\cdot;\mat{\theta})$ 为模型，$\ell$ 为损失函数，$\lambda$ 为正则化系数。目标函数 \eqref{eq:objective} 的物理含义是……

% ============================================================
\section{方法 (Method)}
\subsection{总体框架}
【配一张系统框图。】
\begin{figure}[htbp]
    \centering
    % \includegraphics[width=0.7\textwidth]{framework.png}
    \caption{方法总体框架图}
    \label{fig:framework}
\end{figure}

\subsection{核心模块一}
【数学推导 + 解释。】
\subsection{核心模块二}

% ============================================================
\section{实验 (Experiments)}
\subsection{实验设置}
【数据集（规模/划分）、对比方法、评价指标、实现细节（硬件/超参数/训练轮数）。】
\subsection{对比实验}
\begin{table}[htbp]
    \centering
    \caption{在数据集 X 上的对比结果}
    \label{tab:comparison}
    \begin{tabular}{lcccc}
        \toprule
        \textbf{方法} & \textbf{准确率} & \textbf{F1} & \textbf{参数(百万)} & \textbf{推理耗时(ms)} \\
        \midrule
        Baseline A & 82.1 & 0.80 & 12 & 3.2 \\
        Baseline B & 85.7 & 0.84 & 25 & 5.1 \\
        \textbf{Ours} & \textbf{89.3} & \textbf{0.88} & 18 & 4.0 \\
        \bottomrule
    \end{tabular}
\end{table}
【分析表 \ref{tab:comparison}：为什么更优？分点说明。】
\subsection{消融实验 (Ablation)}
【验证每个组件的贡献。】
\subsection{可视化/案例分析}

% ============================================================
\section{结论 (Conclusion)}
【总结贡献（与引言呼应）+ 不足 + 未来工作。不得引入新内容。】

% ============================================================
%	参考文献（二选一）
%	方案一：natbib + .bib 文件（推荐）
% ============================================================
\bibliographystyle{unsrtnat}
\bibliography{refs}

%	方案二：直接手写（无 .bib 时取消注释下面内容）
% \begin{thebibliography}{99}
% \bibitem{ref1} Author1, Author2. Title of the paper. \emph{Conference/Journal}, 2020.
% \end{thebibliography}

\appendix
\section*{附录}
\subsection{补充证明}
\subsection{补充实验}

\end{document}
